\documentclass[11pt,a4paper]{article}
\usepackage[margin=0.75in]{geometry}
\usepackage{titlesec}
\usepackage{enumitem}
\usepackage{hyperref}

% --- Formatting section headers ---
\titleformat{\section}{\large\bfseries\uppercase}{}{0pt}{}[\titlerule]
\titlespacing*{\section}{0pt}{12pt}{6pt}

\begin{document}
\pagestyle{empty}

% --- HEADER ---
\begin{center}
    {\Huge \textbf{Byungjoo (Daniel) Kang}}  \\[4pt]
    West Lafayette, IN | 765-543-4969 | \href{mailto:kang531@purdue.edu}{kang531@purdue.edu} | LinkedIn | \url{https://daniel-kang-vs.github.io}  
\end{center}

% --- PROFILE ---
\section*{Profile}  
Business Intelligence and Data Analyst specializing in data-driven decision support and Al-assisted analytics, leveraging Python, SQL, Excel, and Power BI to simplify operational and strategic decision-making.   Skilled in end-to-end analytics pipelines, statistical modeling, and executive-level reporting across manufacturing, healthcare, and retail environments.   Strong communicator translating complex analyses into actionable business insights for technical and non-technical stakeholders.  \\[4pt]
\textbf{Skills:} SQL, Python, Power BI, ETL \& Data pipelines (AWS, Azure), Tableau, Excel, PowerPoint  

% --- EDUCATION ---
\section*{Education}  
\noindent
\textbf{Purdue University, Daniels School of Business} \hfill West Lafayette, IN   \\
Master of Science in Business Analytics and Information Management \hfill August 2026   \\
\textit{Certifications:} Google Data Analytics Professional, Microsoft Azure Al Fundamentals   \\[6pt]
\textbf{Purdue University} \hfill West Lafayette, IN   \\
Bachelor of Science in Business Analytics and Information Management, GPA: 3.85/4.0 \hfill December 2024  

% --- PROFESSIONAL EXPERIENCE ---
\section*{Professional Experience}  
\noindent
\textbf{Jambo} \hfill California   \\
\textit{Data Analyst (Remote)} \hfill Jan 2025--August 2025  
\begin{itemize}[leftmargin=*, noitemsep, topsep=0pt]
    \item Engineered analyses to identify key drivers of user churn through funnel and cohort analysis and developed end-to-end ETL analytics pipelines using Python and SQL from behavioral data from 100K+ user events, contributing to a 31\% improvement in retention initiatives.  
    \item Designed and analyzed geo split experiments and multivariate bid tests, applying regression-based lift analysis to validate incrementality and improve customer acquisition efficiency by 27\% while maintaining conversion volume.  
    \item Developed and deployed Al-assisted KPI monitoring dashboards and SQL-based analytics pipelines to automate insight generation and streamline acquisition, engagement, and retention decision-making, reducing reporting time by 60\%.  
\end{itemize}
\vspace{6pt}

\noindent
\textbf{Mayacrew} \hfill Seoul, South Korea   \\
\textit{Data Analyst} \hfill May 2024--August 2024  
\begin{itemize}[leftmargin=*, noitemsep, topsep=0pt]
    \item Developed predictive customer segmentation models using Python and SQL to identify high-value and at-risk user groups across 250K+ customer records, enabling targeted marketing strategies that improved campaign engagement rates by 20\%.  
    \item Built a Python dashboard for B2B engagement KPIs and advised senior leadership on acquisition and retention strategy through clear trade off scenarios, reducing manual reporting time by 40\% and shifting budget toward the highest ROI B2B segments.  
\end{itemize}
\vspace{6pt}

\noindent
\textbf{Pom Education} \hfill Seoul, South Korea   \\
\textit{Deputy Team Leader \& Instructor} \hfill July 2019--Feb 2022  
\begin{itemize}[leftmargin=*, noitemsep, topsep=0pt]
    \item Led student acquisition and consultation efforts by analyzing enrollment funnel and parent decision drivers, contributing to over 1000\% growth in student enrollment in two years.  
    \item Built and maintained a performance tracking system for 1,500+ students using Excel dashboards, analyzing progress, retention, and outcome metrics to identify learning gaps and optimize instructional strategies, improving average student performance by 26\%.  
\end{itemize}

% --- ACADEMIC PROJECTS ---
\section*{Academic Projects}  
\noindent
\textbf{Graduate Data Science Researcher} \hfill West Lafayette, IN   \\
\textit{The Data Mine, Purdue University} \hfill August 2025--Dec 2025  
\begin{itemize}[leftmargin=*, noitemsep, topsep=0pt]
    \item Built an Al-assisted inventory intelligence system integrating predictive reorder-risk modeling and interactive analytics, simplifying SKU prioritization and enabling faster supply chain planning decisions.  
    \item Designed and deployed interactive dashboard and predictive analytics solutions to track reorder likelihood and inventory performance enabling real-time operational monitoring and data-driven supply chain planning.  
\end{itemize}
\vspace{6pt}

\noindent
\textbf{Utilization Analysis and Forecasting Project} \hfill West Lafayette, IN   \\
\textit{Applied Materials Inc.} \hfill January 2024--May 2024  
\begin{itemize}[leftmargin=*, noitemsep, topsep=0pt]
    \item Diagnosed semiconductor equipment utilization patterns by combining time series analysis with macroeconomic indicators, uncovering leading demand signals and delivering insights through dashboards for operations and planning stakeholders.  
    \item Implemented Al-assisted forecasting models combining machine learning and macroeconomic indicators to simplify semiconductor capacity planning and inventory decision-making, improving utilization forecast accuracy by 13\%.  
\end{itemize}

\end{document}